% ============================================================================
% SEC01.TEX - Section 2.1: Sejarah dan Keunggulan Unity
% ============================================================================

\section{Sejarah dan Keunggulan Unity Engine}

\begin{learningoutcomes}
\item Mengetahui sejarah perkembangan Unity Engine
\item Memahami keunggulan Unity dibanding engine lain
\item Mengetahui platform yang didukung Unity
\item Memahami lisensi Unity
\end{learningoutcomes}

\subsection{Sejarah Unity}

Unity Technologies didirikan pada tahun 2004 di Denmark. Unity Engine pertama kali dirilis pada tahun 2005 dengan fokus pada pengembangan game untuk Mac. Sejak itu, Unity telah berkembang menjadi salah satu game engine paling populer di dunia \cite{UnityManual2024}.

\textbf{Timeline Penting}:
\begin{itemize}
    \item \textbf{2005}: Unity 1.0 dirilis (Mac only)
    \item \textbf{2009}: Unity 2.0 (Windows support)
    \item \textbf{2010}: Unity 3.0 (iOS \& Android support)
    \item \textbf{2012}: Unity 4.0 (WebGL, DirectX 11)
    \item \textbf{2017}: Unity 2017 (Timeline, Cinemachine)
    \item \textbf{2020}: Unity 2020 LTS (Stability focus)
    \item \textbf{2023}: Unity 6 (DOTS, Performance)
\end{itemize}

\subsection{Keunggulan Unity Engine}

\begin{enumerate}
\item \textbf{Cross-Platform}: Build sekali, deploy ke banyak platform
\item \textbf{User-Friendly}: Interface yang intuitif untuk pemula
\item \textbf{Asset Store}: Marketplace besar untuk assets dan tools
\item \textbf{Community}: Komunitas besar dan aktif
\item \textbf{Documentation}: Dokumentasi lengkap dan tutorial
\item \textbf{Free Version}: Unity Personal gratis untuk individu dan small teams
\item \textbf{C\# Scripting}: Bahasa pemrograman modern dan powerful
\item \textbf{Visual Editor}: WYSIWYG editor untuk level design
\end{enumerate}

\subsection{Platform yang Didukung}

Unity mendukung deployment ke berbagai platform:
\begin{itemize}
\item Desktop: Windows, Mac, Linux
\item Mobile: iOS, Android
\item Web: WebGL
\item Console: PlayStation, Xbox, Nintendo Switch
\item VR/AR: Oculus, HTC Vive, ARKit, ARCore
\end{itemize}

\subsection{Lisensi Unity}

\begin{itemize}
\item \textbf{Unity Personal}: Gratis untuk individu dan perusahaan dengan revenue < \$100K/tahun
\item \textbf{Unity Pro}: Berbayar dengan fitur tambahan untuk profesional
\item \textbf{Unity Enterprise}: Untuk perusahaan besar dengan kebutuhan khusus
\end{itemize}

\begin{catatan}
Unity Personal sudah cukup untuk pembelajaran dan pengembangan game indie. Upgrade ke Pro hanya diperlukan jika membutuhkan fitur khusus atau untuk komersialisasi skala besar.
\end{catatan}
